% Part: lambda-calculus
% Chapter: introduction
% Section: lambda-definability

\documentclass[../../../include/open-logic-section]{subfiles}

\begin{document}

\olfileid{lam}{int}{rep}
\olsection{\usetoken{S}{lambda definable} حسابی فنکشن}

لامبڈا کیلکولس حساب دا نمونہ کیویں بن سکدا اے؟شروع وچ تاں
ایس گل دا مطلب وی صاف نہیں ہوندا۔قدرتی اعداد اُتے قابلِ حساب ہون
دی گل کرن لئی سانوں ایہناں اعداد دی مناسب نمائندگی لبھنی
پئے گی۔ ہِٹھاں اک ایہو جیہی نمائندگی اے جیہڑی حیرت انگیز
طور اُتے چنگا کم کردی اے۔

\begin{defn}
ہر قدرتی عدد~$n$ لئی \emph{چرچ عدد} $\num{n}$ نوں لامبڈا
ٹرم $\lambd[x][\lambd[y][(x(x(x(\dots
x(y)))))]]$ دے طور اُتے متعین کرو۔ ایتھے کُل $n$ واری $x$ آؤندا اے۔
\end{defn}

ٹرماں $\num{n}$ دا کم عمل دہراؤنا اے۔ انپٹ $f$ اُتے
$\num{n}$ اوہ فنکشن واپس کردا اے جیہڑا $y$ نوں $f^n(y)$
ول بھیجدا اے۔دھیان دیو کہ ہر چرچ عدد نارمل روپ وچ اے۔ہن
اسی دس سکدے آں کہ لامبڈا ٹرم دا قدرتی اعداد اُتے کسے
فنکشن دا «حساب کرنا» کیہہ معنی رکھدا اے۔

\begin{defn}
فرض کرو $f(x_0, \dots, x_{k-1})$ اک $k$-تھاواں والا جزوی
فنکشن اے جیہڑا $\Nat$ توں $\Nat$ ول جاندا اے۔اسی آکھدے آں کہ اک
$\lambd$-ٹرم~$F$ فنکشن~$f$ نوں \emph{!!{lambda define}s}
جے تے صرف جے قدرتی اعداد دے ہر سلسلے $n_0$،
\dots،~$n_{k-1}$ لئی ہِٹھلی تقلیل
\[
F\, \num{n_0}\, \num{n_1} \dots \num{n_{k-1}} \red \num{f(n_0, n_1, \dots,
  n_{k-1})}
\]
اُدوں ہوندی اے جدوں $f(n_0, \dots, n_{k-1})$
متعین ہووے۔ جے اوہ متعین نہ ہووے تاں
$F\, \num{n_0}\, \num{n_1} \dots \num{n_{k-1}}$
دا کوئی نارمل روپ نہ ہووے۔
\end{defn}

\begin{thm}
\ollabel{thm:lambda-def}
اک فنکشن $f$ جزوی قابلِ حساب اے جے تے صرف جے اوہ کسے
لامبڈا ٹرم نال !!{lambda defined} اے۔
\end{thm}

\begin{explain}
ایہ قضیہ کجھ حیران کردا اے۔ حساب دے نمونے دے طور اُتے
لامبڈا کیلکولس کافی سادہ کیلکولس اے؛ایس دے صرف دو عمل
نیں:لامبڈا تجرید تے اطلاق!{} پر ایہناں محدود وسیلیاں
نال وی ہر حسابی طریقۂ کار عمل وچ لیاندا جا سکدا اے۔
\end{explain}

\end{document}
